% !TeX root = ../教材题库_全部条目_转写工作稿.tex
% Source: XJ-B1-C06-PR-P213-Y1382
% 第6章 / 6.1 获取数据的途径与统计的一些基本概念 / 教材 p.213 / PDF 第220页
% OCR未校对：公式与文字均可能有误；未完成数学TeX转写；禁止直接发给学生
% 原题及相邻原解截图（核对用）：../题目原图/XJ-B1-C06-PR-P213-Y1382-01.webp
% 本文件只在转写完成后作为题干使用；OCR 条目尚不是可靠题干。
\begin{bankitem}{p.213 练习 第2题}{第6章 / 6.1 获取数据的途径与统计的一些基本概念 / 教材 p.213 / PDF 第220页}{XJ-B1-C06-PR-P213-Y1382}{OCR未校对：公式与文字均可能有误；未完成数学TeX转写；禁止直接发给学生}
\begin{ocrdraft}
[教材p.213 / PDF第220页]
\par
2．以下调查中，哪些适宜普查，哪些适宜抽样调查？
\par
（D调查某批次汽车的抗撞击能力；
\par
（2）了解某班学生的身高情况；
\par
（3）调查《春节联欢晚会》的收视率；
\par
（4）选出某校短跑最快的学生参加全市比赛．
\end{ocrdraft}
\end{bankitem}
